\documentclass[10pt,a4paper]{amsart}
\usepackage[margin=19mm]{geometry}
\usepackage{amsmath,amssymb,mathtools}
\usepackage[scr=rsfs,cal=euler]{mathalfa}
\usepackage{xfrac}
\usepackage[unicode,hidelinks,bookmarks=false]{hyperref}
\newcommand{\ie}{i.e.\ }
\newcommand{\cf}{cf.\ }
\newcommand{\resp}{respectively\ }
\newcommand{\Subsection}{\subsection}
\providecommand{\nobreakdash}{\nobreak\hbox{-}}
\setlength{\emergencystretch}{2em}
\multlinegap0pt
\usepackage{amssymb}
\usepackage{mathtools}
\usepackage{enumerate}
\usepackage{bm}
\usepackage{tikz}
\newtheorem{theorem}{Theorem}
\newcommand{\C}{\mathbb{C}}
\renewcommand{\Re}{\mathop{\rm Re}}
\DeclareMathOperator{\supp}{supp}
\title{Weighted equilibrium in a field of a uniform charge of an interval}
\author{James Kessinger$^*$, Andrei Mart\'{\i}nez-Finkelshtein}
\date{Source: arXiv:2603.17913v1 (CC BY 4.0)}
\begin{document}
\maketitle

\noindent\emph{Source and licence.} Weighted equilibrium in a field of a uniform charge of an interval. Version arXiv:2603.17913v1 (CC BY 4.0), distributed by arXiv under the Creative Commons Attribution 4.0 International licence, http://creativecommons.org/licenses/by/4.0/, as recorded in the arXiv licence metadata for this exact version and shown on the abstract page https://arxiv.org/abs/2603.17913v1. Full licence notice: \texttt{licenses/arxiv-2603.17913-CC-BY-4.0.txt}.\par\smallskip

\noindent\textit{Editorial scope.} Counted unit: the article's theorem theorem\_case\_-1<=tau<=2/(pi-2) (source0.tex lines 235--259). The article's complete original proof of it is delivered with it (source0.tex lines 260--296, one \texttt{proof} environment of 37 lines). No mathematics is written, repaired or re-derived: the statement and the proof are the article's own lines, in the article's own notation, and no retained line is substituted or re-flowed.  Retained uncounted as the source-local context the proof needs: lines 164--170.  Nothing was omitted from any retained span, so the package carries no omission record.  No retained line contains a citation, so the wrapper carries no bibliography and no bibliography entry.  No retained line contains a figure environment, an image-inclusion command or drawing code, so no image is delivered and the package declares no asset dependency.  Editorial formatting only, in addition: an AMS \texttt{amsart} preamble that restates the article's own declarations and macros used by the retained lines, namely the environments \texttt{theorem} on separate counters, together with the standard packages those lines need.  The retained environments print in this document as Theorem 1 in order of appearance, whereas the article prints the same items with the numbering of its own full document.  The complete native source of the article as distributed in the arXiv e-print for this exact version is bundled unchanged as \texttt{sources/196660/source0.tex}.
\begin{equation}
    \label{equilibrium}
   V^{\mu_\tau}(x)  +\tau V^{\lambda}(x) \begin{cases}
       = \omega_\tau, & x\in \supp (\mu_\tau), \\
    \ge \omega_\tau, & x\in [-1,1].
   \end{cases}
\end{equation}

\begin{theorem}\label{theorem_case_-1<=tau<=2/(pi-2)}
    For $-1\leq\tau\leq\frac{2}{\pi-2}$, $\mu_\tau$ is an absolutely continuous measure on $[-1,1]$ given by
    \begin{equation}\label{mu_when_-1<=tau_<=/(pi-2)}
        d\mu_\tau(x)=\left(\frac{1+\tau}{\pi\sqrt{1-x^2}}-\frac{\tau}{2}\right)\ dx,
    \end{equation}
so that
\begin{equation}\label{Cauchy_Transform_mu_when_-1<=tau_<=/(pi-2)}
        C^{\mu_\tau}(z)= \frac{1+\tau}{\left( z^2-1\right)^{1/2}}-\frac{\tau}{2}\log\left(\frac{z+1}{z-1}\right).
    \end{equation}
Furthermore, for $z\in \C\setminus [-1,1]$,  
\begin{equation*}
\begin{split}
        V^{\mu_\tau}(z)= & -(1+\tau)\log\left|\frac{z+\left( z^2-1\right)^{1/2}}{2}\right|-\tau  \\
        & -\frac{\tau}{2}\left(\operatorname{Re}(z)\log\left|\frac{z-1}{z+1}\right|+\operatorname{Im}(z)\arg\left(\frac{z+1}{z-1}\right)-\log|z^2-1|\right),
\end{split}
    \end{equation*}
and
\begin{equation}\label{Log_Potential_mu_when_-1<=tau_<=/(pi-2)_[-1,1]}
    V^{\mu_\tau}(x)=\frac{\tau}{2} \left[ (1+x)\log(1+x)+ (1-x)\log(1-x) -2 \right] + \omega_\tau, \quad x\in   [-1,1],
\end{equation}
where the equilibrium constant $\omega_\tau$ in \eqref{equilibrium} is  
\begin{equation}\label{equilibrium_constant_-1<=tau<=2/(pi-2)}
    \omega_\tau=(1+\tau)\log(2).
\end{equation}
\end{theorem}

\begin{proof}
    The measure $\mu_\tau$ in \eqref{mu_when_-1<=tau_<=/(pi-2)} is \begin{equation*}\label{Form_of_mu_when_-1<=tau_<=/(pi-2)}
        \mu_\tau=(1+\tau)\eta-\tau\lambda.
    \end{equation*}
    Since $\eta$ and $\lambda$ are probability measures, it follows that \eqref{mu_when_-1<=tau_<=/(pi-2)} is a unit measure. By assumption, $\tau+1>0$, so that the minimum of the density in the right-hand side of \eqref{mu_when_-1<=tau_<=/(pi-2)}, attained for $x=0$, is
    $$
   \frac{1+\tau}{\pi }-\frac{\tau}{2},
    $$
    which is nonnegative due to the upper bound on $\tau$. In consequence, \eqref{mu_when_-1<=tau_<=/(pi-2)} is a positive (and hence a probability) measure on $[-1,1]$.

    Furthermore, for $x\in[-1,1]$,  
    \begin{align*}
      V^{\mu_\tau}(x)+\varphi(x)=   V^{(1+\tau)\eta-\tau\lambda}(x)+\tau V^{\lambda}(x)=(1+\tau)V^\eta(x)=(1+\tau)\log(2).
    \end{align*}
    This proves that \eqref{mu_when_-1<=tau_<=/(pi-2)} is the equilibrium measure and gives us the expression in \eqref{equilibrium_constant_-1<=tau<=2/(pi-2)}. 
    
    Finally, \eqref{Cauchy_Transform_mu_when_-1<=tau_<=/(pi-2)}--\eqref{Log_Potential_mu_when_-1<=tau_<=/(pi-2)_[-1,1]} follow from the known explicit expressions for the Cauchy transform for $\lambda$,
    \begin{equation*} \label{CauchyforLebesgue}
        C^{\lambda}(z)= \frac{1}{2}\log\left(\frac{z+1}{z-1}\right)=\frac{1}{2}\log\left(1+\frac{2}{z-1}\right),
    \end{equation*}
    its primitive,
    \begin{equation} \label{GforLebesgue}
    \begin{split}
        g^{\lambda}(z)&=\int^z C^{\lambda}(y)\, dy= \frac{1}{2} (z+1) \log (z+1)-\frac{1}{2} (z-1) \log (z-1)-1 \\
        &= z \, C^{\lambda}(z)+\frac{1}{2}  \log (z^2-1) -1,
         \end{split}
    \end{equation}
    normalized to have the expansion
    $$
    g^{\lambda}(z)=\log(z)+\mathcal O(z^{-1}), \quad z\to\infty, 
    $$
    and the logarithmic potential of $\lambda$,
    $$
    V^{\lambda}(z)=-\Re g^{\lambda}(z),
    $$
    as well as of the analogous functions of $\eta$.
\end{proof}

\end{document}
